% ID: SEL-D0031-OA-E01
% ESTADO: validada. S032, cotejada con examen y criterios.
\textbf{a)} Multiplicamos el precio de cada artículo por la
cantidad comprada y sumamos los tres importes. Por ejemplo,
la compra de Cati en $C_1$ cuesta $25(2)+20(1)+15(3)=115$.
\[
PQ^t=\begin{pmatrix}115&160\\122&157\end{pmatrix},\qquad
QP^t=(PQ^t)^t=\begin{pmatrix}115&122\\160&157\end{pmatrix}.
\]
En $PQ^t$ las filas son los comercios y las columnas son Cati
y Manuel. En $QP^t$ las filas son las personas y las columnas
los comercios. Cada entrada es el coste de una compra completa.

\textbf{b)} Para Cati es más barato $C_1$, pues $115<122$.
Para Manuel es más barato $C_2$, pues $157<160$.
